\documentclass[12pt,a4paper]{article}
\usepackage[latin1]{inputenc}
\usepackage{float}
\usepackage{amsmath}
\usepackage{amsfonts}
\usepackage{amssymb}
\usepackage{graphicx}
\usepackage[hidelinks]{hyperref}

\author{Giancarlo Danese - 945265}
\author{Davide Savoldelli - } 
\date{A.Y. 2019/2020 - Prof. Di Nitto Elisabetta}


\title{
 \textbf{\Huge{SafeStreets}} \\
 \large Requirements Analysis and Specification Document
}

\begin {document}

 \begin{figure}
  \centering
  \includegraphics[width=1.0\linewidth]{C:/Users/Utente/Documents/logo_poli.jpg}
 \end{figure}

 \maketitle
 \newpage
 \tableofcontents
 \newpage

\section{INTRODUCTION}
\subsection{Purpose} 
This documents represents the Requirements Analysis and Specification Document (RASD) for the SafeStreets software project.
The goal of this project is to deploy a system that allows users to report traffic/street violations to authorities. Users will have the possibility to send pictures of the violation attached with the date/time and position of the vehicle committing the violation. The software focuses particularly on parking violations, including double parking or parking in a disabled-reserved spot.
\subsubsection{Goals}
\begin{itemize}
\item {[G1]}: Allow users to register to the system with their personal data
\item {[G2]}: Allow users to report traffic and parking violations 
\item {[G3]}: Allow end users and authorities to access the information stored
\item {[G4]}: Retrieve the exact position of the place where the violation is being committed
\end{itemize}
\subsection{Scope}
SafeStreets will have an embedded algorithm which will analyze pictures of the vehicle plates sent by the user in order to recognize the vehicle. This information, together with the position of the vehicle and the type of violation that has been committed, will be stored in the software's database. Both end users and mostly authorities will have the chance to mine the information retrieved in the database by highlighting the streets/areas in which most of the violations are committed, the type of vehicles which commit most of the violations and which type of violations occur the most. 		
\subsection{Definitions, Acronyms, Abbreviations} 
\subsection{Revision history} 
\subsection{Reference Documents} 
\subsection{Document Structure} 
	This paper refers to the structure suggested by IEEE for RASD documents, with very slight modifications:
		\begin{enumerate}
			\item \texttt{Introduction}: This first section is a general description of the system's scope and goals. It also includes the revision history of the document and its references, definitions, abbreviations and acronyms used along the paper.
			\item \texttt{Overall Description}: This section includes shared phenomena, requirements and domain assumptions. It also clarifies users' needs.
			\item \texttt{Specific Requirements}: This section is made up of all the requirements needed for the system, both functional and non functional.
			\item \texttt{Formal Analysis Using Alloy}: It includes the Alloy model of the described system.
			\item \texttt{Effort Spent}: In this section you can find information about the hours spent to draft this document.
			\item \texttt{References}: Here are the references about papers/documents used to support this document.
		\end{enumerate}
\section{OVERALL DESCRIPTION}
\subsection{Product perspective: here we include  further details on the shared phenomena and a 
domain model (class diagrams and statecharts)} 
\subsection{Product functions} 
According to the goals defined in the first section, we can list the most important functions the system needs to work correctly:
\begin{itemize}
\item {\texttt{Registration and Login}}: This functionality enables users to register to the system by inserting their personal data, and to login with their chosen username and password. The system must therefore store the data in it's database and retrieve it when queried.  
\item {\texttt{Violation Reporting}}: The system must provide users with an in-app photo sending functionality where users can send pictures of the vehicles plate, insert the type of violation by input and finally write the name of street where the violation occurred.
\item {\texttt{Data Access}}: The software must enable end-users and authorities to access the stored data and mine information from it 
\end{itemize}
\subsection{User characteristics} 
\begin{itemize}
\item {\texttt{Registered User}}: 
\item {\texttt{SafeStreet Administrator}}: 
\item {\texttt{Authorities}}: 
\end{itemize}
\subsection{Assumptions, dependencies and constraints}
 \begin{itemize}
				\item {\textbf[}\textbf{D1}{\textbf]}: The license plate must be clearly visibile in the photo
				\item {\textbf[}\textbf{D2}{\textbf]}: Users' device must support the GPS technology
				\item {\textbf[}\textbf{D3}{\textbf]}: Users' device must be connected to the internet
				\item {\textbf[}\textbf{D4}{\textbf]}: Users' device must have a camera
			\end{itemize}
\section{SPECIFIC REQUIREMENTS: Here we include more details on all aspects in Section 2 if they 
can be useful for the development team}
\subsection{External Interface Requirements} 
\subsubsection{User Interfaces}
\subsubsection{Hardware Interfaces}
\subsubsection{Software Interfaces}
\subsubsection{Communication Interfaces}
\subsection{Functional  Requirements:  Definition  of  use  case  diagrams,  use  cases  and  associated 
sequence/activity diagrams, and mapping on requirements} 
\subsection{Performance Requirements} 
\subsection{Design Constraints}
\subsubsection{Standards compliance}
\subsubsection{Hardware limitations}
\subsubsection{Any other constraint} 
\subsection{Software System Attributes} 
\subsubsection{Reliability}
\subsubsection{Availability}
\subsubsection{Security}
\subsubsection{Maintainability}
\subsubsection{Portability}
\section{FORMAL ANALYSIS USING ALLOY: This section should include a brief presentation of the 
main objectives driving the formal modeling activity, as well as a description of the model 
itself, what can be proved with it, and why what is proved is important given the problem at 
hand. To show  the soundness and correctness of the model,  this section can show some
worlds obtained by running it, and/or the results of the checks performed on meaningful 
assertions}
\section{EFFORT SPENT}
\begin{itemize}
		\item Giancarlo Danese
		\begin{center}
			\begin{tabular}{| c | c | c |}
				\hline
				Day & Subject & Hours \\ \hline
				18/10/2019 & Purpose, Scope & 2 \\
				23/10/2019 & Domain Assumptions, User Characteristics & 1.5 \\
				\hline
			\end{tabular}
		\end{center}

		\item Davide Savoldelli
		\begin{center}
			\begin{tabular}{| c | c | c |}
				\hline
				Day & Subject & Hours \\ \hline
				18/10/2019 & Goals, Product Functions & 1 \\
				23/10/2019 & Document Structure & 2 \\
				\hline
			\end{tabular}
		\end{center}
	\end{itemize}
\section{REFERENCES}


\end {document}